% 右连续单调函数F：x<0时0，0<=x<1时x/2，x>=1时1；跳跃1/2于x=1；区间(0.5,1]的增量F(1)-F(0.5)=1-0.25=0.75。
\begin{tikzpicture}[x=4.2cm,y=2.6cm,
  open/.style={circle,draw=FigureBlue,fill=white,line width=.8pt,inner sep=0pt,minimum size=4.5pt},
  closed/.style={circle,draw=FigureBlue,fill=FigureBlue,line width=.8pt,inner sep=0pt,minimum size=4.5pt},
  note/.style={font=\FigureNoteFont\fontsize{7.5}{9}\selectfont,text=FigureMuted},
  tl/.style={font=\fontsize{7}{8}\selectfont}]
\fill[FigureBlueFill] (0.5,0) rectangle (1,1.08);
\draw[FigureStroke,line width=.7pt,->] (-0.55,0)--(1.62,0) node[right,tl] {$x$};
\draw[FigureStroke,line width=.7pt,->] (0,-0.05)--(0,1.22) node[above,tl] {$F(x)$};
\foreach \y/\l in {0.25/0.25,0.5/0.5,1/1}{\draw[FigureStroke,line width=.5pt] (-0.02,\y)--(0.02,\y);\node[tl,left] at (-0.02,\y) {$\l$};}
\foreach \x/\l in {0.5/0.5,1/1}{\draw[FigureStroke,line width=.5pt] (\x,-0.03)--(\x,0.03);\node[tl,below] at (\x,-0.03) {$\l$};}
\draw[FigureGrid,dashed,line width=.5pt] (0,0.25)--(0.5,0.25);
\draw[FigureGrid,dashed,line width=.5pt] (0,1)--(1,1);
\draw[FigureGrid,dashed,line width=.5pt] (0,0.5)--(1,0.5);
\draw[FigureBlue,line width=1.1pt] (-0.5,0)--(0,0);
\draw[FigureBlue,line width=1.1pt] (0,0)--(1,0.5);
\draw[FigureBlue,line width=1.1pt] (1,1)--(1.55,1);
\node[open] at (1,0.5) {};
\node[closed] at (1,1) {};
\draw[FigureRed,line width=.8pt,<->] (1.08,0.5)--(1.08,1);
\node[note,anchor=west] at (1.1,0.75) {跳跃$\tfrac12$：点质量};
\node[note,anchor=south east] at (0.3,0.16) {斜率$\tfrac12$};
\node[note] at (0.75,1.14) {$\mu_F((0.5,1])=0.75$};
\end{tikzpicture}
